\uuid{Um0a}
\titre{Théorème de transfert}
\auteur{Maxime Nguyen}
\organisation{AMSCC}
\datecreate{2026-11-24}
\niveau{L2}
\module{Probabilité}
\chapitre{Probabilité continue}
\sousChapitre{Espérance et Fonction de répartition}
\difficulte{1}
\theme{espérance, théorème de transfert, loi uniforme}

\contenu{
\begin{question}
Soit $X$ une variable de loi uniforme sur $[0;1]$, notée $X \sim \mathcal{U}([0;1])$ : sa densité vaut $1$ sur $[0;1]$ et $0$ ailleurs. Quelles affirmations sont vraies ?
\begin{answers}
\bad{$\mathbb{E}(X^2) = \mathbb{E}(X)^2 = \frac14$.}
\good{$\mathbb{E}(X^2) = \displaystyle\int_0^1 x^2\,\mathrm{d}x = \frac13$.}
\bad{$\mathbb{E}\left(\frac1X\right) = \dfrac{1}{\mathbb{E}(X)} = 2$.}
\good{$\mathbb{E}\left(\frac1X\right) = \displaystyle\int_0^1 \frac{1}{x}\,\mathrm{d}x = +\infty$.}
\end{answers}
\begin{explanations}
Si $X$ a pour densité $f$, le théorème de
transfert donne $\mathbb{E}(\varphi(X)) = \int \varphi(x) f(x)\,\mathrm{d}x$
pour une fonction $\varphi$ (ici $\varphi(x) = x^2$, puis $\varphi(x) = \frac1x$). En général
$\mathbb{E}(\varphi(X)) \neq \varphi(\mathbb{E}(X))$ ; l'égalité est vraie
quand $\varphi$ est affine.

Contre-sens visé : « l'espérance passe à travers les fonctions ».
\end{explanations}
\end{question}
}
